\chapter{Actor và khái niệm nghiệp vụ}
\label{ch:actor-va-khai-niem}

\section{Hai actor}

Hệ thống có đúng hai actor nghiệp vụ: \textbf{Teacher} và \textbf{Student}.

\subsection{Teacher}

Teacher là người đưa ra các quyết định. Teacher có thể:

\begin{itemize}
  \item Nhập yêu cầu tạo đề.
  \item Xem và chỉnh sửa đề do trợ lí tạo.
  \item Duyệt đề trước khi phát hành.
  \item Phát hành đề cho một hoặc nhiều lớp.
  \item Xem kết quả làm bài của học sinh.
  \item Xem báo cáo học sinh gửi về chỗ trợ lí giải thích chưa rõ.
\end{itemize}

\subsection{Student}

Student là người nhận đề, làm bài và luyện tập. Student có thể:

\begin{itemize}
  \item Xem bài kiểm tra được giao.
  \item Chọn đáp án và nộp bài (giai đoạn 1).
  \item Nghe giải thích lỗi và hỏi lại đến khi hiểu (giai đoạn 2).
  \item Làm câu biến thể của những câu đã sai (giai đoạn 2).
  \item Báo cáo chỗ trợ lí giải thích chưa rõ.
\end{itemize}

\section{Từ điển khái niệm}
\label{sec:tu-dien}

\begin{itemize}
  \item \textbf{Assessment}: Bài kiểm tra, hay bộ câu hỏi được giao cho học sinh.
  \item \textbf{Question}: Câu hỏi trong một assessment.
  \item \textbf{Distractor}: Phương án sai \emph{có chủ đích}, đại diện cho một lỗi hoặc một hiểu nhầm cụ
        thể. Không phải đáp án nhiễu ngẫu nhiên.
  \item \textbf{Submission}: Việc nộp bài, và là mốc kết thúc \emph{giai đoạn 1} --- không phải
        mốc kết thúc bài kiểm tra.
  \item \textbf{Attempt}: Cả một lần làm bài của một học sinh trên một đề, trải qua \textbf{cả hai giai
        đoạn}. Mỗi học sinh có \emph{nhiều nhất} một Attempt cho mỗi đề --- em chưa vào làm thì
        không có Attempt nào. Nó xong khi mọi câu sai đã chốt, hoặc khi hạn giai đoạn 2 đi qua.
  \item \textbf{Feedback}: Phản hồi / nhận xét của trợ lí cho những thắc mắc của học sinh.
  \item \textbf{Misconception}: Hiểu nhầm khái niệm, hoặc lỗi sai có tính lặp lại.
  \item \textbf{Mastery}: Mức độ thành thạo của học sinh với một mục tiêu học tập. Khái niệm này vẫn còn
        trong từ điển nhưng \textbf{không còn quyết định gì} --- điều kiện dừng của vòng chữa bài
        là số vòng, không phải mastery.
  \item \textbf{Giai đoạn 1}: Phần làm bài và nộp.
  \item \textbf{Giai đoạn 2}: Phần chữa bài.
  \item \textbf{Lượt}: Một lượt cho phép học sinh làm lại các câu sai của giai đoạn trước / lượt trước.
  \item \textbf{Câu biến thể}: Câu sinh từ \emph{câu học sinh làm sai ở giai đoạn 1}: giữ cấu trúc và lỗi cần
        kiểm, chỉ đổi dữ kiện.
  \item \textbf{Lời giải}: Cách làm đi kèm mỗi câu hỏi, bắt buộc có nhiều hơn một cách. Trợ lí
        dựa vào lời giải để giải thích cho học sinh.
  \item \textbf{Class}: Lớp học có sẵn danh sách học sinh. Giáo viên tạo lớp. Phát hành đề phải
        chọn một lớp đã tạo từ trước.
  \item \textbf{Question Bank}: Ngân hàng đề.
  \item \textbf{Document}: Tài liệu PDF của giáo viên, nằm trong kho lưu trữ riêng của từng tài
        khoản. Cung cấp kiến thức và giới hạn phạm vi ra đề.
  \item \textbf{Scope}: Phạm vi ra đề, chương và khoảng trang trong một \emph{Document}.
  \item \textbf{Service}: Một đơn vị chạy độc lập trong hệ thống. Có năm: \texttt{frontend},
        \texttt{backend}, \texttt{agent}, \texttt{document}, \texttt{ingest}.
  \item \textbf{Job} và \textbf{Task}: Hai thứ khác nhau. Một \emph{task} là \emph{tên} của một loại
        việc. Một \emph{job} là một lần chạy của task ấy, có id riêng và có thể thất bại riêng.
  \item \textbf{Queue}: Nơi một service đẩy job vào và một service khác lấy ra làm. Job nằm lại
        trên queue đến khi có ai lấy, nên tắt máy giữa chừng không mất việc.
  \item \textbf{Worker}: Tiến trình ngồi nghe một queue, lấy job xuống và chạy. Nó không nhận lời
        gọi từ ngoài; thứ duy nhất nó đọc là queue.
  \item \textbf{Pub/sub}: Phát một tiếng lên một \emph{channel}: ai đang nghe thì nhận, ai không
        nghe thì mất. Khác queue ở chỗ nó \textbf{không giữ lịch sử} và không chở nội dung --- nó
        chỉ nói \emph{có cái gì vừa đổi}, còn muốn biết đổi gì thì phải đi hỏi lại.
  \item \textbf{SSE} và \textbf{Polling}: Hai cách để màn hình biết có thứ mới. Với \emph{SSE},
        \texttt{backend} giữ một kết nối mở và \emph{hích} xuống mỗi khi có gì đổi. Với \emph{polling},
        màn hình tự hỏi lại theo một chu kì. Hệ thống dùng \emph{SSE}; \emph{polling} là phương án
        nó \textbf{không} dùng, và có trong từ điển để gọi tên thứ đã bị bỏ.
  \item \textbf{Gateway}: Cửa duy nhất đi vào một vùng. \texttt{backend} là gateway của cả hệ thống:
        không màn hình nào gọi thẳng \texttt{agent}, \texttt{document} hay \texttt{ingest}.
  \item \textbf{Store}: Ba kho của hệ thống là ba loại khác nhau, không nên gọi chung một chữ:
        \emph{database} giữ dữ liệu có cấu trúc, \emph{object store} (một \emph{bucket}) giữ byte
        của tệp, và \emph{queue} giữ job đang chờ.
  \item \textbf{Draft}: Đề đang soạn, chưa duyệt. Câu hỏi ở trạng thái này sửa và xoá được tự do.
  \item \textbf{Approve}: Chốt rằng đề đã dùng được. Đây là một \emph{trạng thái}, không phải một
        thao tác: nó có cả cạnh ngược (\emph{unapprove}), và chỉ đề đã approve mới phát hành được.
  \item \textbf{Publish} và \textbf{Publication}: \emph{Publish} là phát hành một đề cho một lớp. Một
        \emph{Publication} là bản ghi sinh ra từ việc đó: một lớp đã chọn, cộng năm thứ giáo viên
        đặt cho lớp đó.
        Cùng một đề phát hành cho hai lớp là hai \emph{Publication}.
  \item \textbf{Mark} và \textbf{Score}: \emph{Mark} là điểm của \emph{một câu}; \emph{score} là điểm
        của \emph{cả bài}. Tiếng Việt gọi cả hai là ``điểm'', nên báo cáo giữ hai chữ tiếng Anh ở
        những chỗ cần phân biệt.
  \item \textbf{opens\_at} / \textbf{closes\_at} / \textbf{remediation\_deadline}: Ba mốc thời
        gian của một \emph{Publication}. Hai mốc đầu mở và đóng giai đoạn 1; mốc thứ ba là hạn
        của giai đoạn 2. Chữ ``hạn'' một mình không nói rõ là hạn nào.
  \item \textbf{Tool}: Một việc có tên mà \texttt{agent} được phép \emph{đề xuất} chạy, kèm tham
        số. Danh sách tool là một \textbf{ràng buộc}, không phải một tiện ích: duyệt đề và phát
        hành đề cố ý không có tool nào, nên \texttt{agent} không có đường nào để đề xuất chúng.
\end{itemize}
